\chapter{The problem: finding what does not belong}\label{ch:problem}

\begin{onething}
We never teach the computer what an anomaly looks like. We teach it what \emph{normal} looks like
in a scene, give every pixel a number that says how badly it fits that normal, and hand the
highest-scoring pixels to a person.
\end{onething}

This chapter has no heavy mathematics. Its job is to fix the question that the rest of the book
answers, to show with arithmetic why the obvious answer fails, and to set up the notation that
every later chapter uses. Read it slowly anyway: almost every design decision in the forty
chapters that follow is a consequence of the three facts in \cref{sec:problem:classifier}.

\section{The situation}\label{sec:problem:situation}

A satellite flies over a patch of ground and records an image. Somebody at a mining ministry, an
environmental agency or an insurance company wants to know one thing:
\emph{is there anything unusual down there?}

Unusual might mean an illegal mine, a leaking pipeline, a chemical dumped in a field, an
unusually hot piece of machinery, a vehicle parked where no vehicle should be, or a mineral
outcrop that could be worth money. The software this book describes is called
\textbf{Allotrope}\index{Allotrope}. It is a web application: a user uploads a satellite file,
runs a chain of processing steps on it, picks a cut-off on a map of scores, and downloads a
bundle naming the most unusual places with their coordinates. \Cref{ch:product} describes the
product; this chapter is about the question it answers.

Before any mathematics, three words.

\begin{definition}[Pixel]\index{pixel}
A \textbf{pixel} is one cell of the grid an image is made of. In a satellite image a pixel
stands for a patch of ground; for the sensors in this book that patch is about 30 metres on a
side (\cref{ch:light}). We name a pixel by its row $i$ and column $j$, both counted from $0$.
\end{definition}

\begin{definition}[Scene]\index{scene}
A \textbf{scene} is one complete satellite acquisition: one patch of the Earth, at one moment,
as one file (or folder of files).
\end{definition}

\begin{definition}[Anomaly]\index{anomaly}
An \textbf{anomaly} is a pixel (or a small group of pixels) that is unlike the rest of its
scene. The word says nothing about \emph{what} the pixel is. In this codebase ``anomaly'',
``outlier'' and ``novelty'' are used interchangeably; the code says ``anomaly'' throughout.
\end{definition}

\section{Why a classifier will not work}\label{sec:problem:classifier}

If you have met machine learning before, your instinct is a \textbf{classifier}%
\index{classifier}: a program that takes a pixel and answers with one of a fixed list of
labels. You would collect thousands of pictures of illegal mines, label them ``illegal mine'',
\emph{train} the classifier on them (adjust its internal numbers until its answers on the
labelled examples are right), and run it on new imagery.

That approach is standard, and here it fails for three separate reasons. Each one alone would
sink it, and each closes a different escape route, so it is worth keeping all three.

\subsection{Reason 1: there are no labels}

A \textbf{label}\index{label} is the correct answer attached to a training example. To label a
satellite pixel ``illegal mine'' somebody has to know what is on the ground under it, which in
practice means somebody visited the spot. Nobody has done that for thousands of anomalies across
every terrain and season. The number of labelled anomalies available is, for practical purposes,
zero. A classifier with no labelled examples has nothing to learn from.

\subsection{Reason 2: you do not know what you are looking for}

A classifier can only answer with labels it was given in advance. Train it on ``illegal mine''
and ``oil spill'' and it will find illegal mines and oil spills. The valuable find is usually
the one nobody anticipated: a chemical nobody had seen before, a structure nobody expected. A
classifier is blind to unnamed classes \emph{by construction}, and the unnamed classes are the
point of surveying.

\subsection{Reason 3: the interesting thing is vanishingly rare}

Anomalies are a tiny fraction of the pixels. That rarity breaks the ordinary way of measuring
whether a classifier works.

\begin{definition}[Accuracy]\index{accuracy}
The \textbf{accuracy} of a yes/no rule over $N$ pixels is the fraction of pixels it answers
correctly:
\begin{equation}\label{eq:problem:accuracy}
\text{accuracy} = \frac{\text{number of correct answers}}{N}.
\end{equation}
\end{definition}

\begin{worked}[label={we:problem:accuracy}]{The ``always normal'' rule}
A scene has $N = 20{,}000{,}000$ pixels, of which $40$ are anomalous. Consider the rule that
ignores its input and answers ``normal'' for every pixel. Work out its accuracy on paper before
reading on.

The rule is right on every normal pixel and wrong on every anomalous one. Count the normal
pixels:
\begin{align*}
\text{normal pixels} &= 20{,}000{,}000 - 40\\
                     &= 19{,}999{,}960.
\end{align*}
Every one of those is a correct answer, so put that count into \cref{eq:problem:accuracy}:
\begin{align*}
\text{accuracy} &= \frac{19{,}999{,}960}{20{,}000{,}000}\\
                &= 0.999998\\
                &= 99.9998\%.
\end{align*}
The rule finds nothing and scores $99.9998\%$. The fraction of anomalous pixels, often called
the \textbf{base rate}\index{base rate}, is
\begin{align*}
\frac{40}{20{,}000{,}000} &= \frac{4}{2{,}000{,}000}\\
                          &= 0.000002 = 2\times 10^{-6}.
\end{align*}
Whenever the base rate is this small, accuracy is close to $1$ for every rule that says
``normal'' almost all the time, useful or not.
\end{worked}

A classifier trained on such data learns the same lesson as the ``always normal'' rule, because
that rule is nearly always right. Accuracy is therefore useless here; \cref{ch:measures} builds
the measures used instead.

\begin{keypoint}
No labels, an open-ended target and extreme rarity together rule out supervised
classification; any one of them alone would be enough.
\end{keypoint}

\section{The reframe: model normality, score the residue}\label{sec:problem:reframe}

Stop asking ``is this pixel a mine?'', which needs labels. Ask instead:
\begin{center}
\emph{Given everything else in this scene, how surprising is this pixel?}
\end{center}
That question needs no labels at all. It needs only a description of what normal looks like
here, and every scene carries its own definition of normal in its own pixels. It also catches
things nobody anticipated, because ``unlike everything else here'' does not require knowing what
the thing is. The sentence to remember from this chapter is:

\begin{center}
\textbf{We never model anomalies. We model normality, and call the residue anomaly.}
\end{center}

Everything in the rest of this book is machinery for two jobs: building a model of normality,
and measuring the residue.

\begin{definition}[Anomaly score and score map]\index{anomaly score}\index{score map}
An \textbf{anomaly score} is a function $s$ that gives each pixel $(i,j)$ of a scene a number
$s(i,j) \ge 0$, where larger means ``fits the model of normal worse''. Laid out on the pixel
grid, the scores form a \textbf{score map}: one number per pixel, the same height and width as
the image.
\end{definition}

The score map is the product's central output. It is not a yes/no answer and it is not a list of
objects. A person looks at it, chooses a cut-off, and decides what the highest-scoring pixels
are. The system's output is \textbf{a shortlist for an analyst, not a verdict}.

\begin{analogy}
A proofreader working in a language they know well has no list of the typos in this book. What
they have is a deep sense of what correct words look like. At ``the cat sat on the mta'' their
eye snags: they did not recognise a known error, they \emph{failed to recognise a normal word}.
That snag is the anomaly score. The proofreader can tell you \emph{where} something is off without
being able to say what the author meant, and in the same way the system can say \emph{where}
without saying \emph{what}. You give up naming the thing in exchange for not needing to be told
what to look for. The analogy stops working at one point: a proofreader corrects the typo, while
the system never proposes what the pixel ``should'' have been as a finding, only how far it is
from what it expected.
\end{analogy}

\section{Two ways to model normality}\label{sec:problem:families}

This codebase has exactly two families of method. Both ship in the product, and a user can run
both on the same scene and compare.

\subsection{Family 1: classical detectors, fitted to each scene}

A \textbf{classical detector}\index{classical detector} computes statistics of the scene in
front of it: the average pixel, and how much pixels typically vary. It then measures how far each
pixel is from that average, in units of the typical variation. Nothing is trained in advance and
nothing is stored between scenes; the statistics are computed fresh from each scene's own pixels.

The simplest version of this idea uses one number per pixel, which is what a thermal sensor
gives (\cref{ch:light}). Write $x$ for a pixel's value, $\mu$ (``mu'') for the average over the
scene, and $\sigma$ (``sigma'') for the typical spread. The distance from the average in units of
spread is
\begin{equation}\label{eq:problem:z}
z = \frac{x - \mu}{\sigma},
\end{equation}
and a score can be its size, $s = \abs{z}$. Read \cref{eq:problem:z} in words: subtract the
average to get ``how far above or below average'', then divide by the spread to measure that
distance in standard units.

\begin{definition}[Mean and standard deviation]\index{mean}\index{standard deviation}
For $n$ numbers $x_1, \dots, x_n$, the \textbf{mean} and the \textbf{standard deviation} are
\[ \mu = \frac{1}{n}\sum_{k=1}^{n} x_k, \qquad \sigma = \sqrt{\frac{1}{n}\sum_{k=1}^{n} (x_k - \mu)^2}. \]
For the mean, add the numbers up and divide by how many there are. For the standard deviation,
take each number's distance from the mean, square it so negatives do not cancel positives, average the squares, and take the square
root to return to the original units.
\end{definition}

\begin{worked}[label={we:problem:zscore}]{A one-band scene with one hot pixel}
A tiny thermal scene has five pixels with surface temperatures, in degrees Celsius,
\[ 20,\ 21,\ 19,\ 20,\ 30. \]
Do each line on paper. First the mean:
\begin{align*}
\text{sum} &= 20 + 21 + 19 + 20 + 30 = 110,\\
\mu &= 110 / 5 = 22.
\end{align*}
Next each difference from the mean, and its square:
\begin{align*}
20 - 22 &= -2, & (-2)^2 &= 4,\\
21 - 22 &= -1, & (-1)^2 &= 1,\\
19 - 22 &= -3, & (-3)^2 &= 9,\\
20 - 22 &= -2, & (-2)^2 &= 4,\\
30 - 22 &= 8,  & 8^2 &= 64.
\end{align*}
Average the squares, then take the square root:
\begin{align*}
\text{sum of squares} &= 4 + 1 + 9 + 4 + 64 = 82,\\
\sigma^2 &= 82/5 = 16.4,\\
\sigma &= \sqrt{16.4} \approx 4.050.
\end{align*}
Now the score of each pixel, $s = \abs{x-\mu}/\sigma$:
\begin{align*}
s(20) &= 2/4.050 \approx 0.49, & s(21) &= 1/4.050 \approx 0.25, & s(19) &= 3/4.050 \approx 0.74,\\
s(30) &= 8/4.050 \approx 1.98.
\end{align*}
The hot pixel scores about four times higher than any other. Notice also that the hot pixel
pulled the mean up from $20$ to $22$ and inflated the spread; a detector that computes its
statistics from the very pixels it scores lets an anomaly partly hide itself. \Cref{ch:rx} and
\cref{ch:localrx} return to this.
\end{worked}

The classical method in this book is called \textbf{RX}\index{RX detector}
\citep{reed1990adaptive}. It is \cref{eq:problem:z} for many numbers per pixel at once, using a
distance (the Mahalanobis distance, \cref{ch:rx}) that accounts for the fact that some directions
in the data vary more than others. The version shipped for hyperspectral data is
\textbf{MNF-RX}\index{MNF-RX}, which first transforms the data into a smaller, better-behaved
space \citep{green1988transformation}.

\begin{drift}
\emph{History, dated 2026-05-11.} Plain RX on the 165-band hyperspectral data was deliberately
dropped on that date: with that many bands the covariance matrix it inverts is nearly singular
and the distances blew up to around $10^{11}$. MNF-RX replaced it. Do not reintroduce plain RX on
hyperspectral data; it has been tried. (Source: the project instructions \file{CLAUDE.md} and the
Indradhanu course, part 1.) \Cref{ch:rx,ch:localrx} explain why the inversion fails.
\end{drift}

\subsection{Family 2: foundation models, trained once to redraw imagery}

The second family uses a \textbf{neural network}\index{neural network}: a function built from
many simple layers whose behaviour is set by a large list of adjustable numbers called
\textbf{parameters}\index{parameter} (\cref{ch:shapes} defines layers precisely). The network is
\textbf{trained}, once and offline, on tens of thousands of ordinary satellite images, to do a
task that needs no labels: hide part of an image and redraw the hidden part from its
surroundings. Training produces a file of learned parameters (``weights''); running the model on
a new scene loads that file.

In this repository such a network is called a \textbf{foundation model}\index{foundation model}.
Careful with that phrase: in the wider industry it usually means an enormous language model with
hundreds of billions of parameters. Here it means something much smaller and more specific, a
network trained to reconstruct imagery whose reconstruction error is used as the anomaly score.
The hyperspectral model in this book, \textbf{Indradhanu}\index{Indradhanu} (Sanskrit for
``rainbow''), had $5{,}507{,}354$\tagM{indradhanu workings/05-where-it-lives.md} parameters in the
version the courses describe; \cref{ch:params} counts them by hand. Its thermal sibling is
\textbf{Chakshu}\index{Chakshu} (``the eye'').

The rule that turns a reconstruction into a score is ``reconstruct, then subtract''
(\cref{fig:problem:pipeline}). Where the model redraws the scene well, the difference between
what was observed and what was redrawn is small. Where the model, trained only on ordinary
ground, meets something it has never seen, it paints ordinary background there instead, and the
difference is large.

\coursefig[\linewidth]{iw-01-pipeline}{Reconstruct, then subtract. Blocks of the observed cube
are hidden and the model, trained only on ordinary scenes, redraws them from their surroundings.
Subtracting the reconstruction from the observation, band by band, gives a per-pixel error map;
where the redrawing is most wrong (amber) is where something unusual is. The amber square stands
for an object the model has never seen. (Course figure, Indradhanu part 1.)}{fig:problem:pipeline}

\begin{worked}[label={we:problem:residual}]{A reconstruction residual on three numbers}
A pixel is described by three numbers (three ``bands'', \cref{ch:light}): observed
$\vect{x} = (0.10,\ 0.20,\ 0.30)$. The model's redrawing of it is
$\xhat = (0.11,\ 0.18,\ 0.30)$. These values are illustrative\tagA. One simple score is the mean
absolute difference over the three numbers. Do each step.

Subtract band by band:
\begin{align*}
0.10 - 0.11 &= -0.01,\\
0.20 - 0.18 &= 0.02,\\
0.30 - 0.30 &= 0.00.
\end{align*}
Take absolute values (drop the signs, because over- and under-shooting are equally bad):
\[ 0.01,\quad 0.02,\quad 0.00. \]
Add them and divide by $3$:
\begin{align*}
0.01 + 0.02 + 0.00 &= 0.03,\\
s &= 0.03/3 = 0.01.
\end{align*}
If a second pixel had the same observation but the model redrew it as $(0.10,\ 0.10,\ 0.10)$,
the differences would be $0.00$, $0.10$, $0.20$, their sum $0.30$, and the score
$0.30/3 = 0.10$: ten times larger. That pixel is the one the model could not explain.
\Cref{ch:recon} explains why a pixel must never be allowed to help reconstruct itself, and
\cref{ch:losses,ch:residualmaha} replace this simple average with better scores.
\end{worked}

\begin{table}[htbp]\centering
\caption{The two families of anomaly model in this codebase.}\label{tab:problem:families}
\begin{tabular}{@{}lll@{}}\toprule
 & Classical detectors & Foundation models\\\midrule
Model of normal & scene mean and spread & a trained network \\
Trained? & no, fitted per scene & yes, once, offline \\
Stored file of weights & none & a checkpoint (\texttt{.pt}) \\
Score & distance from the mean & reconstruction error \\
Examples here & MNF-RX & Indradhanu, Chakshu \\\bottomrule
\end{tabular}
\end{table}

\section{From a score map to a shortlist: thresholds are percentiles}\label{sec:problem:threshold}

A score map ranks pixels; someone still has to decide where to cut. The obvious rule is an
absolute one: ``flag every pixel scoring above $0.05$''. It fails, and the reason is worth seeing
with numbers.

\begin{worked}[label={we:problem:fixedcut}]{Why a fixed cut-off fails}
Two scenes are scored by the same model. The numbers are illustrative\tagA, of the size the
Indradhanu course quotes.
\begin{itemize}
\item Scene A, a calm desert: typical score $0.002$, highest score $0.04$.
\item Scene B, a hazy coastal city: typical score $0.08$.
\end{itemize}
Apply the rule ``flag if $s > 0.05$''. In scene A the highest score is $0.04$ and $0.04 < 0.05$,
so nothing is flagged, including the real anomaly. In scene B the typical pixel already scores
$0.08 > 0.05$, so most of the image is flagged. The typical score differs between the scenes by a
factor of
\[ 0.08 / 0.002 = 40, \]
so no single number can serve both.
\end{worked}

The cure is to make the cut relative to the scene in front of you.

\begin{definition}[Percentile threshold]\index{percentile}\index{threshold}
Sort a scene's $N$ valid scores from largest to smallest. The \textbf{top-$q$ threshold}, for a
fraction $q$ such as $0.001$ (``the top $0.1\%$''), keeps the $k$ highest-scoring pixels, where
$k$ is $qN$ rounded to a whole number. Equivalently, the cut-off value is the
$(1-q)$ \textbf{percentile} of the scores: the value below which a fraction $1-q$ of them lie.
\end{definition}

\begin{worked}[label={we:problem:percentile}]{A top-20\% cut on ten scores}
Ten pixels have scores
\[ 0.3,\ 0.1,\ 0.9,\ 0.2,\ 0.4,\ 0.2,\ 0.7,\ 0.1,\ 0.3,\ 0.2. \]
Keep the top $q = 0.2$. Do it on paper.

How many pixels to keep:
\[ k = q N = 0.2 \times 10 = 2. \]
Sort the scores from largest to smallest:
\[ 0.9,\ 0.7,\ 0.4,\ 0.3,\ 0.3,\ 0.2,\ 0.2,\ 0.2,\ 0.1,\ 0.1. \]
Keep the first two: the pixels scoring $0.9$ and $0.7$. Any cut-off strictly between $0.4$ and
$0.7$ gives the same set. If every score in the scene were multiplied by $10$, the sorted order
would not change and the same two pixels would be kept; that indifference to the overall scale is
exactly what a fixed cut-off lacked.

For a real scene of $1{,}210 \times 1{,}219$ pixels, all valid, the top $0.1\%$ is
\begin{align*}
N &= 1{,}210 \times 1{,}219 = 1{,}474{,}990,\\
k &= 0.001 \times 1{,}474{,}990 = 1{,}474.99 \approx 1{,}475 \text{ pixels.}
\end{align*}
\end{worked}

In the product a person makes this choice by dragging a slider over the score map, in a separate
step that deliberately pauses for them (\cref{ch:product}). Separating ``score'' from ``cut'' also
means one expensive scoring run can be tried at five different thresholds.

\begin{keypoint}
Thresholds in this system are relative to the scene (percentiles), never absolute numbers.
\end{keypoint}

\section{Why satellites make this harder}\label{sec:problem:satellites}

Two complications follow you through the whole book.

\subsection{Scale}

{\sloppy A single scene from the PRISMA satellite (\cref{ch:sensors}) is about $1{,}210$ by $1{,}219$
pixels with $239$ measurements per pixel. Both figures are quoted in module docstrings, the
first in the PRISMA reflectance transformer in \file{app/utils/data_transformations/}
\allowbreak\tagM{prs_l2d_dn_to_surface_reflectance_transformer.py} and the second in the PRISMA
dataset builder in \file{app/utils/dataset_builder/} \allowbreak\tagM{prisma_dataset_builder.py}.
Count the numbers:\par}
\begin{align*}
1{,}210 \times 1{,}219 &= 1{,}474{,}990 \text{ pixels},\\
1{,}474{,}990 \times 239 &= 352{,}522{,}610 \text{ numbers}.
\end{align*}
Stored as 4-byte floating-point values that is
\[ 352{,}522{,}610 \times 4 = 1{,}410{,}090{,}440 \text{ bytes} \approx 1.41\ \text{GB} \]
for one copy of one scene, which is why the code worries about how many copies exist at once.

\subsection{Not every pixel is real}

Sensors have dead detectors. Clouds block the ground. A satellite sweeps a strip, so a
rectangular file has ragged, empty edges where the strip did not reach. A large fraction of any
scene is data you must actively ignore.

\begin{definition}[Validity mask]\index{validity mask}
A \textbf{validity mask} is a companion array of the same grid as the image holding $1$ where
the pixel is a genuine measurement and $0$ where it is not.
\end{definition}

Forgetting the mask is the single most common source of false detections in this codebase.
Invalid pixels typically hold $0$ or a special ``fill'' value, and a detector looking for unusual
values finds a large block of identical extreme values and ranks it top.

\begin{worked}[label={we:problem:unmasked}]{An unmasked zero becomes the top anomaly}
Take the five temperatures of \cref{we:problem:zscore} and add a sixth pixel that is off the
edge of the swath, stored as $0$. If the detector forgets the mask, the scene is
\[ 20,\ 21,\ 19,\ 20,\ 30,\ 0. \]
Mean:
\begin{align*}
\text{sum} &= 20 + 21 + 19 + 20 + 30 + 0 = 110,\\
\mu &= 110/6 \approx 18.33.
\end{align*}
Squared differences from $18.33$:
\begin{align*}
(20 - 18.33)^2 &= 1.67^2 \approx 2.79,\\
(21 - 18.33)^2 &= 2.67^2 \approx 7.13,\\
(19 - 18.33)^2 &= 0.67^2 \approx 0.45,\\
(20 - 18.33)^2 &= 1.67^2 \approx 2.79,\\
(30 - 18.33)^2 &= 11.67^2 \approx 136.19,\\
(0 - 18.33)^2 &= (-18.33)^2 \approx 335.99.
\end{align*}
Sum, average, root:
\begin{align*}
2.79 + 7.13 + 0.45 + 2.79 + 136.19 + 335.99 &= 485.34,\\
\sigma^2 &= 485.34/6 \approx 80.89,\\
\sigma &\approx 8.99.
\end{align*}
Scores of the two extreme pixels:
\begin{align*}
s(30) &= 11.67/8.99 \approx 1.30,\\
s(0) &= 18.33/8.99 \approx 2.04.
\end{align*}
The fake zero now outranks the genuinely hot pixel, and its presence has halved the hot pixel's
score (from $1.98$ to $1.30$). With the mask applied, the zero is excluded from both the
statistics and the ranking, and \cref{we:problem:zscore} is recovered.
\end{worked}

The validity mask appears in \cref{ch:light} as a ``validity cube'', becomes binary per pixel in
\cref{ch:bandpipeline}, and decides which tokens a network may see in \cref{ch:masking}.

\section{Notation for the rest of the book}\label{sec:problem:notation}

The symbols below are fixed from here on. Later chapters add to the list but do not change it.

\begin{table}[htbp]\centering
\caption{Notation used throughout the book.}\label{tab:problem:notation}
\begin{tabular}{@{}lp{0.68\linewidth}@{}}\toprule
Symbol & Meaning \\\midrule
$C$, $H$, $W$ & number of bands, rows (height) and columns (width) of an image\\
$c$ & a band index, $0 \le c \le C-1$\\
$(i,j)$ & a pixel at row $i$, column $j$, both counted from $0$\\
$\mat{X}$ & a whole image (``cube''), an array of shape \shape{(C, H, W)}\\
$x_c(i,j)$ & the value of band $c$ at pixel $(i,j)$\\
$\vect{x}$, $\vect{x}_{ij}$ & one pixel's $C$ values as a column vector (its spectrum, \cref{ch:light})\\
$\xhat$ & a model's reconstruction of $\vect{x}$\\
$v(i,j)$ & validity: $1$ for a real measurement, $0$ otherwise\\
$s(i,j)$ & anomaly score of pixel $(i,j)$, larger means more unusual\\
$\mu$, $\sigma$ & a mean and a standard deviation; $\muv$, $\Cov$ their many-band versions\\
$\lambda$ & wavelength, in nanometres (\cref{ch:light})\\\bottomrule
\end{tabular}
\end{table}

Indices start at $0$, as they do in Python, because the book quotes Python code constantly. A
cube with $C = 165$ bands has bands $0$ to $164$.

\begin{confusions}
\confusion{So the model is trained to detect anomalies?}{No. A foundation model is trained to
redraw ordinary imagery; it has no concept of an anomaly. Detection happens afterwards, by
subtracting. A classical detector is not trained at all.}
\confusion{Does it need examples of anomalies to work?}{No, and that is the point. Neither family
is ever shown an anomaly on purpose. Labelled anomalies are used only to \emph{measure} a
detector afterwards (\cref{ch:measures}), never to build one.}
\confusion{Is a high score a guarantee of something interesting?}{No. A high score means ``the
model could not explain this''. Sensor glitches, cloud edges and unusual-but-boring terrain also
score high. Turning a high score into a finding is the job of later steps and of a person.}
\confusion{Anomaly detection is outlier detection, then?}{Related, but where the comparison is made
matters. A pixel can hold a perfectly ordinary temperature for the scene as a whole and still be
anomalous because everything right around it is 20 degrees colder. Several detectors score a pixel
against its neighbours rather than the whole scene (\cref{ch:localrx}).}
\confusion{Why not pick one good absolute threshold by experiment?}{Because score magnitudes
change with the scene by an order of magnitude or more (\cref{we:problem:fixedcut}). A cut that
works on one scene flags nothing or everything on another. Percentiles do not care about the
overall scale.}
\confusion{Does ``foundation model'' here mean a large language model?}{No. Here it means a small
network (a few million parameters) trained to reconstruct imagery, whose reconstruction error is
the score.}
\end{confusions}

\begin{checkyourself}
\item\label{ex:problem:1} Give the three reasons a supervised classifier is the wrong tool here,
one sentence each, and say which escape route each one closes.
\item\label{ex:problem:2} A scene has $1{,}000{,}000$ pixels and $10$ are anomalous. Compute, step
by step, the accuracy of the rule ``always answer normal'' and the base rate.
\item\label{ex:problem:3} A one-band scene has the four values $10$, $12$, $14$, $20$. Compute the
mean, the standard deviation (dividing by $n$), and the score $\abs{x-\mu}/\sigma$ of the pixel
reading $20$. Show every line.
\item\label{ex:problem:4} A scene has $2{,}000{,}000$ valid pixels. How many pixels does a top
$0.05\%$ cut keep? If every score is doubled, does the kept set change? Why?
\item\label{ex:problem:5} A three-band pixel is observed as
$(0.05,\ 0.25,\ 0.40)$; the model reconstructs it as $(0.07,\ 0.22,\ 0.40)$. Compute the mean absolute difference. Then explain in two
sentences why a pixel stored as $0$ because it lies off the swath must be masked out before any
such score is ranked.
\end{checkyourself}

\nextchapter{ch:light}
